\begin{table}[htbp]
\caption{Cumulative flood damage and casualty indicators for the ten regencies with the highest Flood Severity Index (2016--2025)}
\label{tab:fsi_components}
\centering
\footnotesize
\begin{tabular*}{\textwidth}{@{\extracolsep{\fill}}llrrrrrrr@{}}
\toprule
\textbf{Regency / City} & \textbf{Province} & \textbf{\shortstack{Flood\\Events}} & \textbf{Deaths} & \textbf{Missing} & \textbf{Injured} & \textbf{\shortstack{Houses\\Flooded}} & \textbf{\shortstack{Houses\\Damaged}} & \textbf{\shortstack{Facilities\\Damaged}} \\
\midrule
Aceh Utara & Aceh & 48 & 240 & 6 & 920 & 79,796 & 47,208 & 583 \\
Jayapura & Papua & 20 & 106 & 82 & 983 & 4,964 & 1,148 & 220 \\
Agam & Sumatera Barat & 34 & 221 & 39 & 281 & 3,094 & 1,976 & 203 \\
Padang Pariaman & Sumatera Barat & 46 & 53 & 6 & 524 & 26,503 & 3,853 & 795 \\
Sukabumi & Jawa Barat & 87 & 9 & 1 & 10,382 & 5,311 & 5,499 & 78 \\
Luwu Utara & Sulawesi Selatan & 101 & 39 & 9 & 107 & 52,766 & 4,177 & 134 \\
Banjar & Kalimantan Selatan & 38 & 11 & 0 & 871 & 223,701 & 19,817 & 1,360 \\
Bogor & Jawa Barat & 277 & 20 & 1 & 25 & 52,247 & 3,499 & 187 \\
Tapanuli Selatan & Sumatera Utara & 36 & 96 & 6 & 253 & 3,603 & 2,890 & 110 \\
Garut & Jawa Barat & 77 & 36 & 21 & 10 & 11,240 & 4,915 & 431 \\
\bottomrule
\end{tabular*}
\vspace{4pt}
\begin{minipage}{\textwidth}
\footnotesize \textit{Note:} Data compiled from the National Disaster Management Authority (BNPB DIBI, 2016--2025). Values represent cumulative raw counts.
\end{minipage}
\end{table}
